% Source-only Beamer presentation. Simulation result figures are placeholders.
% After running the Python code, copy the files from
% 01_Report_Figures/<timestamp> into figures/simulation/.
\documentclass[aspectratio=169,10pt]{beamer}
\usetheme{Madrid}
\setbeamertemplate{navigation symbols}{}
\setbeamertemplate{footline}[frame number]

\usepackage{amsmath,amssymb,graphicx,booktabs,xcolor,ragged2e}
\usepackage[most]{tcolorbox}
\usepackage{adjustbox}

\definecolor{DeepBlue}{RGB}{18,45,74}
\definecolor{SoftBlue}{RGB}{44,112,173}
\definecolor{Accent}{RGB}{0,135,160}
\definecolor{LightGray}{RGB}{245,247,250}
\definecolor{SoftRed}{RGB}{180,35,35}
\definecolor{DarkGreen}{RGB}{35,105,75}

\setbeamercolor{title}{fg=white,bg=DeepBlue}
\setbeamercolor{frametitle}{fg=white,bg=DeepBlue}
\setbeamercolor{structure}{fg=SoftBlue}
\setbeamercolor{block title}{fg=white,bg=SoftBlue}
\setbeamercolor{block body}{fg=black,bg=LightGray}
\setbeamerfont{title}{size=\Large,series=\bfseries}
\setbeamerfont{frametitle}{size=\large,series=\bfseries}
\setbeamerfont{block body}{size=\footnotesize}
\setbeamertemplate{itemize items}[circle]

\newcommand{\kw}[1]{\textcolor{Accent}{\textbf{#1}}}
\newcommand{\gap}[1]{\textcolor{SoftRed}{\textbf{#1}}}
\newcommand{\paperfig}[2][]{\includegraphics[#1]{figures/paper/#2}}
\newcommand{\simfig}[3][]{%
\IfFileExists{figures/simulation/#2}{\includegraphics[#1]{figures/simulation/#2}}{%
\begin{tcolorbox}[colback=white,colframe=SoftBlue,title={#3},arc=1mm,boxrule=0.8pt,width=\linewidth]
\centering
\vspace{0.28cm}
\textbf{Insert simulation figure after running the Python code.}\\[0.2cm]
Expected filename:\\[-0.05cm]
\texttt{\detokenize{#2}}\\[0.15cm]
\scriptsize Copy it from \texttt{Simulation\_Output/<timestamp>/For\_Slides} into \texttt{figures/simulation/}.
\vspace{0.28cm}
\end{tcolorbox}}}


% Direct paper evidence boxes: used instead of unclear screenshot highlighting.
\newtcolorbox{quoteevidence}[2][]{
  colback=white,
  colframe=#2,
  boxrule=0.9pt,
  arc=1.2mm,
  left=1.5mm,right=1.5mm,top=1mm,bottom=1mm,
  fonttitle=\bfseries\scriptsize,
  title={#1}
}
\newcommand{\paperquote}[2]{%
\begin{quoteevidence}[Direct phrase from the paper]{#1}
\small ``#2''
\end{quoteevidence}}

\newcommand{\partslide}[2]{%
{\setbeamercolor{background canvas}{bg=DeepBlue}%
\begin{frame}[plain]
\vfill
\centering
{\color{white}\Huge\textbf{#1}}\\[0.45cm]
{\color{white}\Large #2}
\vfill
\end{frame}}}

\title[LSTM System Identification]{LSTM-Based Dynamic System Identification}
\subtitle{Wang (2017), Ogunmolu/Gans (2016), and Actuator Current-to-Displacement/Force Modeling}
\author{Hussein Zolfaghari}
\date{}

\begin{document}

\begin{frame}[plain,noframenumbering]
\titlepage
\end{frame}


\begin{frame}[t]{Presentation roadmap}
\vspace{0.03cm}
\scriptsize
\begin{columns}[T,onlytextwidth]

\begin{column}{0.32\textwidth}
\begin{tcolorbox}[
enhanced,
colback=LightGray,
colframe=DeepBlue,
colbacktitle=DeepBlue,
coltitle=white,
title=\textbf{Part I: Wang (2017)},
fonttitle=\scriptsize\bfseries,
boxrule=0.8pt,
arc=2mm,
left=1.8mm,right=1.8mm,top=1.4mm,bottom=1.4mm,
height=0.58\textheight
]
\centering
{\large\textbf{Convex-based LSTM}}
\vspace{0.08cm}
\RaggedRight
\begin{itemize}\itemsep2pt
\item Problem and research gap
\item Series--parallel identification
\item Multiple LSTMs with convex weights
\item Simulated nonlinear-system results
\end{itemize}
\end{tcolorbox}
\end{column}

\begin{column}{0.32\textwidth}
\begin{tcolorbox}[
enhanced,
colback=LightGray,
colframe=SoftBlue,
colbacktitle=SoftBlue,
coltitle=white,
title=\textbf{Part II: Ogunmolu/Gans},
fonttitle=\scriptsize\bfseries,
boxrule=0.8pt,
arc=2mm,
left=1.8mm,right=1.8mm,top=1.4mm,bottom=1.4mm,
height=0.58\textheight
]
\centering
{\large\textbf{Deep Dynamic NNs}}
\vspace{0.08cm}
\RaggedRight
\begin{itemize}\itemsep2pt
\item Nonlinear system identification
\item Hammerstein-type modeling
\item RNN, LSTM, FastLSTM, GRU
\item Soft-robot and GlassFurnace data
\end{itemize}
\end{tcolorbox}
\end{column}

\begin{column}{0.32\textwidth}
\begin{tcolorbox}[
enhanced,
colback=LightGray,
colframe=Accent,
colbacktitle=Accent,
coltitle=white,
title=\textbf{Part III: My simulation},
fonttitle=\scriptsize\bfseries,
boxrule=0.8pt,
arc=2mm,
left=1.8mm,right=1.8mm,top=1.4mm,bottom=1.4mm,
height=0.58\textheight
]
\centering
{\large\textbf{Actuator Modeling}}
\vspace{0.08cm}
\RaggedRight
\begin{itemize}\itemsep2pt
\item Input: coil current
\item Outputs: displacement and force
\item LSTM next-step prediction
\item Full and zoomed tracking results
\end{itemize}
\end{tcolorbox}
\end{column}

\end{columns}
\end{frame}


\partslide{Part I}{Wang (2017): LSTM Neural Networks for Dynamic System Identification}

\begin{frame}{Wang: paper objective, gap, and contribution}
\begin{columns}[T,onlytextwidth]
\begin{column}{0.36\textwidth}
\begin{block}{Main point}
The paper studies nonlinear dynamic system identification using LSTM neural networks.
\end{block}
\begin{block}{Research gap}
The paper argues that LSTM had been successful in machine learning, but its use for control-oriented dynamic system identification had not been sufficiently addressed.
\end{block}
\begin{block}{Novelty}
Instead of using only one LSTM, Wang combines several LSTM models with adaptive convex weights to make the identification error converge faster.
\end{block}
\end{column}
\begin{column}{0.60\textwidth}
\centering
\paperfig[width=\linewidth,height=0.80\textheight,keepaspectratio]{paper_title_abstract.png}
\end{column}
\end{columns}
\end{frame}

\begin{frame}{Wang: analytical model question}
\begin{columns}[T,onlytextwidth]
\begin{column}{0.38\textwidth}
\begin{block}{Question}
Did the paper produce a final classical linear model of the form
\[
\dot{x}=Ax+Bu?
\]
\end{block}
\begin{block}{Answer}
\gap{No.} The paper starts from a general nonlinear dynamic representation and then formulates nonlinear discrete-time identification using input-output histories.
\end{block}
\end{column}
\begin{column}{0.58\textwidth}
\centering
\paperfig[width=\linewidth,height=0.80\textheight,keepaspectratio]{general_dynamic_system_overview_infographic.png}
\end{column}
\end{columns}
\end{frame}



\begin{frame}{Wang methodology: series-parallel identification}
\begingroup
\scriptsize
\setlength{\abovedisplayskip}{3pt}
\setlength{\belowdisplayskip}{3pt}

\begin{columns}[T,onlytextwidth]

\begin{column}{0.34\textwidth}

\begin{block}{Clear meaning}

\textbf{Parallel identification} = uses past inputs and past predicted outputs:
\[
[u_{\mathrm{past}},\hat{y}_{\mathrm{past}}]\rightarrow \hat{y}(k)
\]

\textbf{Series path} = uses past outputs:
\[
y(k-1),\ldots,y(k-n)
\]

\textbf{Series--parallel} = uses both histories to predict the next output:
\[
\hat{y}(k)=f(u\text{-history},y\text{-history})
\]
\end{block}

\vspace{-0.12cm}

\begin{block}{Link to my simulation}
My LSTM follows the same idea. It uses the last 120 samples of current, displacement, and force to predict the next displacement and force:
\[
\text{past }[I,x,F]\rightarrow \text{next }[x,F].
\]
\end{block}

\end{column}

\begin{column}{0.62\textwidth}
\centering
\paperfig[width=\linewidth,height=0.78\textheight,keepaspectratio]{series_parallel_identification_model_diagram.png}
\end{column}

\end{columns}
\endgroup
\end{frame}

\begin{frame}{Wang methodology: convex-based LSTM cluster}
\begin{columns}[T,onlytextwidth]

\begin{column}{0.35\textwidth}

\begin{block}{Clear meaning}
\footnotesize
Wang does not rely on only one LSTM model.

Instead, several LSTM models work in parallel:
\[
\text{LSTM}_1,\ \text{LSTM}_2,\ \ldots,\ \text{LSTM}_n
\]

Each one gives its own prediction:
\[
\hat{y}_1,\hat{y}_2,\ldots,\hat{y}_n.
\]
\end{block}

\begin{block}{Final prediction}
\footnotesize
The predictions are combined using convex weights:
\[
\hat{y}
=
\sum_{i=1}^{n}\alpha_i\hat{y}_i,
\qquad
\alpha_i\ge0,\quad
\sum_{i=1}^{n}\alpha_i=1.
\]

This means the final output is a weighted average of all LSTM predictions.
\end{block}

\end{column}

\begin{column}{0.61\textwidth}
\centering
\paperfig[width=\linewidth,height=0.78\textheight,keepaspectratio]{convex_based_lstm_ensemble_diagram.png}
\end{column}

\end{columns}
\end{frame}

\begin{frame}{Wang: data usage and training style}
\begin{columns}[T,onlytextwidth]

\begin{column}{0.35\textwidth}

\begin{block}{Simulation data source}
\footnotesize
The paper uses a \textbf{simulated nonlinear SISO system}:
\[
u(k) \rightarrow y(k).
\]
The output is generated from the nonlinear difference equation shown on the right, not from experimental actuator measurements.
\end{block}

\begin{block}{Identification setup}
\footnotesize
\begin{itemize}\itemsep2pt
\item The input signal is analytically defined.
\item The model is updated over time using the identification error
\[
e(k)=y(k)-\hat{y}(k).
\]
\item No explicit train/validation/test split, batch size, or epoch count is reported.
\end{itemize}
\end{block}

\end{column}

\begin{column}{0.61\textwidth}
\centering
\paperfig[width=\linewidth,height=0.9\textheight,keepaspectratio]{wang_simulation1_data_training_highlighted.png}

\vspace{0.03cm}
{\scriptsize\color{SoftRed}
}
\end{column}

\end{columns}
\end{frame}

\begin{frame}{Wang: reported simulation result}
\begin{columns}[T,onlytextwidth]

\begin{column}{0.33\textwidth}
\begin{block}{Comparison}
\footnotesize
The paper compares the identification error:
\[
e(k)=y(k)-\hat{y}(k)
\]
for two methods:
\begin{itemize}\itemsep2pt
\item conventional RNN
\item convex-based LSTM
\end{itemize}
\end{block}

\begin{block}{Main observation}
\footnotesize
The convex-based LSTM error becomes small faster and remains smoother than the RNN error.
\end{block}

\begin{block}{Interpretation}
\footnotesize
The result supports the paper's claim that the convex LSTM structure improves convergence speed.
\end{block}
\end{column}

\begin{column}{0.63\textwidth}
\centering
\paperfig[width=\linewidth,height=0.78\textheight,keepaspectratio]{comparison_of_rnn_and_lstm_errors.png}
\end{column}

\end{columns}
\end{frame}





\partslide{Part II}{Ogunmolu/Gans (2016): Deep Dynamic Neural Networks for Nonlinear System Identification}

\begin{frame}{Ogunmolu/Gans: problem, gap, and contribution}
\begin{columns}[T,onlytextwidth]
\begin{column}{0.55\textwidth}
\centering
\paperfig[width=\linewidth,height=0.34\textheight,keepaspectratio]{ogun_excerpt_goal.png}
\vspace{0.10cm}
\paperfig[width=\linewidth,height=0.31\textheight,keepaspectratio]{ogun_excerpt_gap_contribution.png}
\end{column}
\begin{column}{0.41\textwidth}
\begin{block}{Problem definition}
\small
Use deep neural-network models to identify nonlinear dynamic systems from sequential input--output data.
\end{block}
\begin{block}{Research gap}
\small
Highly nonlinear dynamics can be too complicated for a simple closed-form analytical model.
\end{block}
\begin{block}{Main contribution}
\small
The paper evaluates NN-based Hammerstein and recurrent/LSTM-type dynamic models on SISO and MIMO system-identification datasets.
\end{block}
\end{column}
\end{columns}
\vspace{-0.04cm}
{\scriptsize\color{gray}Only the paper lines supporting the goal, gap, and contribution are boxed in red.}
\end{frame}

\begin{frame}{Ogunmolu/Gans: analytical model question}
\begin{columns}[T,onlytextwidth]
\begin{column}{0.48\textwidth}
\begin{block}{Question}
Did the paper end with a final analytical model such as
\[
\dot{x}=Ax+Bu?
\]
\end{block}
\begin{block}{Answer}
\gap{No.} The paper presents data-driven nonlinear input--output identification models. It does not identify one final constant state-space pair $(A,B)$.
\end{block}
\end{column}
\begin{column}{0.48\textwidth}
\paperfig[width=\linewidth,height=0.46\textheight,keepaspectratio]{ogun_excerpt_gap_contribution.png}
\begin{block}{How to say it in the presentation}
\small
The authors use neural networks because the nonlinear dynamics are difficult to write as closed-form equations. Therefore, the result is a learned dynamic estimator, not an explicit linear state-space law.
\end{block}
\end{column}
\end{columns}
\end{frame}

\begin{frame}{Ogunmolu/Gans methodology: Hammerstein model}
\begin{columns}[T,onlytextwidth]
\begin{column}{0.47\textwidth}
\centering
\paperfig[width=\linewidth,height=0.36\textheight,keepaspectratio]{ogun_excerpt_hammerstein_fig5.png}
\vspace{0.10cm}
\paperfig[width=\linewidth,height=0.36\textheight,keepaspectratio]{ogun_excerpt_hammerstein_text.png}
\end{column}
\begin{column}{0.49\textwidth}
\begin{block}{Main modeling idea}
\small
The Hammerstein structure separates the model into:
\begin{itemize}\itemsep2pt
\item a nonlinear neural-network block, $g(\cdot)$;
\item a linear dynamic block, $G(z^{-1})$;
\item an output prediction $y(n)$.
\end{itemize}
\end{block}
\begin{block}{Why this is relevant}
\small
This is a dynamic input--output model. It supports the same general direction as my simulation: learning actuator dynamics directly from measured time-series data.
\end{block}
\end{column}
\end{columns}
\end{frame}

\begin{frame}{Ogunmolu/Gans methodology: LSTM architecture}
\begin{columns}[T,onlytextwidth]
\begin{column}{0.56\textwidth}
\centering
\paperfig[width=\linewidth,height=0.76\textheight,keepaspectratio]{ogun_excerpt_lstm_architecture.png}
\end{column}
\begin{column}{0.40\textwidth}
\begin{block}{Reported architecture}
\small
The paper replaces the recurrent nonlinear element with LSTM modules:
\[
1\rightarrow1,\quad 1\rightarrow10,\quad 10\rightarrow100,\quad 100\rightarrow1.
\]
\end{block}
\begin{block}{Presentation point}
\small
The important idea is not a static input--output fit. The paper uses recurrent dynamic neural networks to include time-history effects.
\end{block}
\end{column}
\end{columns}
\end{frame}

\begin{frame}{Ogunmolu/Gans: data and training/test usage}
\begin{columns}[T,onlytextwidth]
\begin{column}{0.57\textwidth}
\centering
\paperfig[width=\linewidth,height=0.70\textheight,keepaspectratio]{ogun_excerpt_data_training.png}
\end{column}
\begin{column}{0.39\textwidth}
\begin{block}{Soft-robot dataset}
\small
\begin{itemize}\itemsep2pt
\item SISO experiment.
\item Input: valve current in mA.
\item Output: head motion in mm.
\item 10,070 offline input--output samples.
\end{itemize}
\end{block}
\begin{block}{Training/test protocol}
\small
\begin{itemize}\itemsep2pt
\item 60\% training, 40\% testing.
\item Mini-batch size: 100 samples.
\item Training: 50 epochs.
\item No validation split is reported.
\end{itemize}
\end{block}
\begin{block}{Benchmark dataset}
\small
GlassFurnace: 3 inputs and 6 outputs.
\end{block}
\end{column}
\end{columns}
\end{frame}

\begin{frame}{Ogunmolu/Gans: reported experimental results}
\begin{columns}[T,onlytextwidth]
\begin{column}{0.56\textwidth}
\centering
\paperfig[width=\linewidth,height=0.34\textheight,keepaspectratio]{ogun_excerpt_results_softrobot_table.png}
\vspace{0.10cm}
\paperfig[width=\linewidth,height=0.32\textheight,keepaspectratio]{ogun_excerpt_results_glassfurnace_table.png}
\end{column}
\begin{column}{0.40\textwidth}
\begin{block}{Soft-robot result}
\small
The reported fit values are very high for all tested models. LSTM gives high fit, but it is not the lowest-MSE model in this table.
\end{block}
\begin{block}{GlassFurnace result}
\small
The benchmark result shows that model architecture matters. Vanilla LSTM performs poorly, while FastLSTM improves the result.
\end{block}
\begin{block}{Main lesson}
\small
Dynamic neural networks can model nonlinear systems, but architecture and data usage strongly affect the final accuracy.
\end{block}
\end{column}
\end{columns}
\end{frame}




% -------------------------------------------------------------------------
% Three clear comparison slides: Wang vs. Ogunmolu/Gans and my simulation
% -------------------------------------------------------------------------

\begin{frame}{Paper 1 vs. Paper 2: main difference}
\scriptsize
\setlength{\abovedisplayskip}{3pt}
\setlength{\belowdisplayskip}{3pt}

\begin{columns}[T,onlytextwidth]

\begin{column}{0.485\textwidth}
\begin{tcolorbox}[
enhanced,
colback=LightGray,
colframe=DeepBlue,
colbacktitle=DeepBlue,
coltitle=white,
title=\textbf{Paper 1: Wang (2017)},
fonttitle=\small\bfseries,
boxrule=0.8pt,
arc=2mm,
left=2mm,right=2mm,top=1.5mm,bottom=1.5mm,
height=0.58\textheight
]
\textbf{Main focus}\par
Build a new LSTM-based \kw{identifier} for nonlinear dynamic systems.

\vspace{0.10cm}
\textbf{Main method}\par
Use measured input and output histories:
\[
[u_{\rm past},y_{\rm past}]\rightarrow\hat y(k).
\]
Combine several LSTMs by convex weighting:
\[
\hat y=\sum_i \alpha_i\hat y_i.
\]

\vspace{0.08cm}
\textbf{Main result}\par
Faster and smoother identification-error convergence than conventional RNN.

\vspace{0.08cm}
\textbf{Data type}\par
Mostly simulated nonlinear SISO systems.
\end{tcolorbox}
\end{column}

\begin{column}{0.485\textwidth}
\begin{tcolorbox}[
enhanced,
colback=LightGray,
colframe=SoftBlue,
colbacktitle=SoftBlue,
coltitle=white,
title=\textbf{Paper 2: Ogunmolu/Gans (2016)},
fonttitle=\small\bfseries,
boxrule=0.8pt,
arc=2mm,
left=2mm,right=2mm,top=1.5mm,bottom=1.5mm,
height=0.58\textheight
]
\textbf{Main focus}\par
Compare \kw{deep dynamic neural-network architectures} for system identification.

\vspace{0.10cm}
\textbf{Main method}\par
Evaluate several dynamic models:
\[
\text{MLP},\ \text{RNN},\ \text{LSTM},\ \text{FastLSTM},\ \text{GRU}.
\]
Use recurrent and Hammerstein-type models, not Wang's explicit series--parallel structure.

\vspace{0.08cm}
\textbf{Main result}\par
Report fit percentage, MSE, training time, and model comparison.

\vspace{0.08cm}
\textbf{Data type}\par
Real soft-robot and GlassFurnace data.
\end{tcolorbox}
\end{column}

\end{columns}

\vspace{0.08cm}
\begin{tcolorbox}[
enhanced,
colback=LightGray,
colframe=Accent,
colbacktitle=Accent,
coltitle=white,
title=\textbf{One-sentence takeaway},
fonttitle=\small\bfseries,
boxrule=0.7pt,
arc=1.5mm,
left=2mm,right=2mm,top=1mm,bottom=1mm,
height=0.15\textheight
]
\textbf{Wang} focuses on \kw{how to build the identifier}; \textbf{Ogunmolu/Gans} focuses on \kw{which dynamic neural architecture works better on real data}.
\end{tcolorbox}
\end{frame}


\begin{frame}{What I used from each paper in my actuator simulation}
\scriptsize
\setlength{\abovedisplayskip}{3pt}
\setlength{\belowdisplayskip}{3pt}

\begin{columns}[T,onlytextwidth]

\begin{column}{0.32\textwidth}
\begin{tcolorbox}[
enhanced,
colback=LightGray,
colframe=DeepBlue,
colbacktitle=DeepBlue,
coltitle=white,
title=\textbf{From Wang},
fonttitle=\small\bfseries,
boxrule=0.8pt,
arc=2mm,
left=2mm,right=2mm,top=1.5mm,bottom=1.5mm,
height=0.66\textheight
]
\begin{itemize}\itemsep3pt
\item Dynamic input--output view, not static curve fitting.
\item Series--parallel training with measured histories:
\[
[I,x,F]_{\rm past}\rightarrow[\hat x,\hat F]_{\rm next}.
\]
\item Parallel/free-running rollout for stronger testing:
\[
[I,\hat x,\hat F]_{\rm past}\rightarrow[\hat x,\hat F]_{\rm next}.
\]
\item Error and convergence tracking.
\end{itemize}
\end{tcolorbox}
\end{column}

\begin{column}{0.32\textwidth}
\begin{tcolorbox}[
enhanced,
colback=LightGray,
colframe=SoftBlue,
colbacktitle=SoftBlue,
coltitle=white,
title=\textbf{From Ogunmolu/Gans},
fonttitle=\small\bfseries,
boxrule=0.8pt,
arc=2mm,
left=2mm,right=2mm,top=1.5mm,bottom=1.5mm,
height=0.66\textheight
]
\begin{itemize}\itemsep3pt
\item Apply dynamic neural networks to real experimental data.
\item Compare model behavior instead of assuming one model is best.
\item Report quantitative metrics:
\[
\text{MSE},\ \text{RMSE},\ R^2,\ \text{fit}.
\]
\item Show measured--predicted and error plots.
\end{itemize}
\end{tcolorbox}
\end{column}

\begin{column}{0.32\textwidth}
\begin{tcolorbox}[
enhanced,
colback=LightGray,
colframe=Accent,
colbacktitle=Accent,
coltitle=white,
title=\textbf{My final setup},
fonttitle=\small\bfseries,
boxrule=0.8pt,
arc=2mm,
left=2mm,right=2mm,top=1.5mm,bottom=1.5mm,
height=0.66\textheight
]
\begin{itemize}\itemsep4pt
\item \textbf{Input:} $I(k)$, coil current.
\item \textbf{Outputs:} $x(k)$ and $F(k)$.
\item \textbf{Window:}
\[
[I,x,F]_{\text{past }120}.
\]
\item \textbf{Target:}
\[
[x,F]_{\rm next}.
\]
\item \textbf{Final test:} unseen chirp experiment.
\end{itemize}
\end{tcolorbox}
\end{column}

\end{columns}
\end{frame}


\begin{frame}{Why I implemented these ideas}
\scriptsize
\begin{columns}[T,onlytextwidth]

\begin{column}{0.485\textwidth}
\begin{tcolorbox}[
enhanced,
colback=LightGray,
colframe=DeepBlue,
colbacktitle=DeepBlue,
coltitle=white,
title=\textbf{Why Wang's ideas are useful},
fonttitle=\small\bfseries,
boxrule=0.8pt,
arc=2mm,
left=2mm,right=2mm,top=1.5mm,bottom=1.5mm,
height=0.52\textheight
]
\begin{itemize}\itemsep4pt
\item My actuator is dynamic; the current effect depends on recent history.
\item Series--parallel training is stable because measured past outputs are supplied.
\item Free-running rollout is a stronger test because the model feeds back its own predictions.
\item Error curves show whether the model learns dynamics, not only points.
\end{itemize}
\end{tcolorbox}
\end{column}

\begin{column}{0.485\textwidth}
\begin{tcolorbox}[
enhanced,
colback=LightGray,
colframe=SoftBlue,
colbacktitle=SoftBlue,
coltitle=white,
title=\textbf{Why Ogunmolu/Gans's ideas are useful},
fonttitle=\small\bfseries,
boxrule=0.8pt,
arc=2mm,
left=2mm,right=2mm,top=1.5mm,bottom=1.5mm,
height=0.52\textheight
]
\begin{itemize}\itemsep4pt
\item My data are real experimental time series, not only a simulated equation.
\item Comparing dynamic neural models helps justify the selected architecture.
\item Metrics make the results measurable and not only visual.
\item Prediction and error plots make tracking quality easy to explain.
\end{itemize}
\end{tcolorbox}
\end{column}

\end{columns}

\vspace{0.10cm}
\begin{tcolorbox}[
enhanced,
colback=LightGray,
colframe=Accent,
colbacktitle=Accent,
coltitle=white,
title=\textbf{How this guides my simulation},
fonttitle=\small\bfseries,
boxrule=0.8pt,
arc=2mm,
left=2mm,right=2mm,top=1.4mm,bottom=1.4mm,
height=0.22\textheight
]
I combine the two lessons: \textbf{Wang} guides the identification structure and rollout tests, while \textbf{Ogunmolu/Gans} guides practical evaluation on experimental data. Therefore, my simulation reports \kw{one-step prediction}, \kw{closed-loop unseen-chirp simulation}, and quantitative metrics.

\vspace{0.03cm}
{\tiny\color{gray}\textbf{Note:} I use ideas from both papers, but I do not claim that Ogunmolu/Gans used Wang's exact series--parallel implementation.}
\end{tcolorbox}
\end{frame}



\partslide{Part III}{My Actuator Simulation: Series, Series--Parallel, and Parallel Results}

\begin{frame}{My simulation: data usage and unseen test}
\scriptsize
\begin{columns}[T,onlytextwidth]
\begin{column}{0.34\textwidth}
\begin{block}{Development data}
\begin{itemize}\itemsep2pt
\item Four chirp experiments:
\[
67,\ 87,\ 107,\ 127\ \mathrm{mA}.
\]
\item Random non-overlapping blocks are used for training and validation.
\item This helps cover different chirp-frequency regions.
\end{itemize}
\end{block}
\begin{block}{Final test}
\[
147\ \mathrm{mA}\rightarrow\text{unseen test only}.
\]
No normalization, training, or validation uses this record.
\end{block}
\end{column}
\begin{column}{0.62\textwidth}
\centering
\simfig[width=\linewidth,height=0.78\textheight,keepaspectratio]{02_data_usage_summary.png}{Training, validation, and test data usage}
\end{column}
\end{columns}
\end{frame}

\begin{frame}{My simulation: three implemented identification structures}
\scriptsize
\setlength{\abovedisplayskip}{3pt}
\setlength{\belowdisplayskip}{3pt}
\begin{columns}[T,onlytextwidth]

\begin{column}{0.32\textwidth}
\begin{tcolorbox}[
enhanced,
colback=LightGray,
colframe=DeepBlue,
colbacktitle=DeepBlue,
coltitle=white,
title=\textbf{Series},
fonttitle=\small\bfseries,
boxrule=0.8pt,
arc=2mm,
left=2mm,right=2mm,top=1.4mm,bottom=1.4mm,
height=0.56\textheight
]
\textbf{Input-only model}
\[
[I,\Delta I,I_{DC}]_{\rm past}
\rightarrow
[\hat x,\hat F]_{\rm next}
\]
No output history is supplied.
\end{tcolorbox}
\end{column}

\begin{column}{0.32\textwidth}
\begin{tcolorbox}[
enhanced,
colback=LightGray,
colframe=SoftBlue,
colbacktitle=SoftBlue,
coltitle=white,
title=\textbf{Series--parallel},
fonttitle=\small\bfseries,
boxrule=0.8pt,
arc=2mm,
left=2mm,right=2mm,top=1.4mm,bottom=1.4mm,
height=0.56\textheight
]
\textbf{Measured-feedback model}
\[
[I,\Delta I,I_{DC},x,F]_{\rm past}
\rightarrow
[\hat x,\hat F]_{\rm next}
\]
Measured output history is supplied.
\end{tcolorbox}
\end{column}

\begin{column}{0.32\textwidth}
\begin{tcolorbox}[
enhanced,
colback=LightGray,
colframe=Accent,
colbacktitle=Accent,
coltitle=white,
title=\textbf{Parallel},
fonttitle=\small\bfseries,
boxrule=0.8pt,
arc=2mm,
left=2mm,right=2mm,top=1.4mm,bottom=1.4mm,
height=0.56\textheight
]
\textbf{Free-running model}
\[
[I,\Delta I,I_{DC},\hat x,\hat F]_{\rm past}
\rightarrow
[\hat x,\hat F]_{\rm next}
\]
Predictions are fed back recursively.
\end{tcolorbox}
\end{column}

\end{columns}
\vspace{0.05cm}
\begin{block}{Comparison question}
Which structure best predicts displacement and Lorentz force on the unseen 147 mA chirp?
\end{block}
\end{frame}

\begin{frame}{My simulation: training idea}
\scriptsize
\begin{columns}[T,onlytextwidth]
\begin{column}{0.48\textwidth}
\begin{block}{Architecture}
The known current features first pass through a static nonlinear block:
\[
[I,\Delta I,I_{DC}]\rightarrow g_{\theta}(\cdot).
\]
Then an LSTM models the dynamic time-history relation.
\end{block}
\begin{block}{Training stages}
\begin{itemize}\itemsep2pt
\item Series: trained from input history only.
\item Series--parallel: trained with measured output history.
\item Parallel: starts from series--parallel weights and is fine-tuned recursively.
\end{itemize}
\end{block}
\end{column}
\begin{column}{0.48\textwidth}
\begin{block}{Why random blocks?}
The chirp frequency changes over time. Random non-overlapping blocks help training and validation include different frequency regions without leaking nearly identical LSTM windows.
\end{block}
\begin{block}{How to interpret accuracy}
Series--parallel is usually easier. Parallel is harder because prediction errors are reused as future inputs.
\end{block}
\end{column}
\end{columns}
\end{frame}

\begin{frame}{Simulation result: convergence of all structures}
\scriptsize
\begin{columns}[T,onlytextwidth]
\begin{column}{0.31\textwidth}
\begin{block}{What to check}
\begin{itemize}\itemsep2pt
\item Training loss trend.
\item Validation loss trend.
\item Differences between the three structures.
\end{itemize}
\end{block}
\begin{block}{Important note}
Parallel loss is a harder recursive objective, so it should not be compared blindly with one-step loss.
\end{block}
\end{column}
\begin{column}{0.65\textwidth}
\centering
\simfig[width=\linewidth,height=0.78\textheight,keepaspectratio]{01_all_structures_cost.png}{Training and validation cost}
\end{column}
\end{columns}
\end{frame}


\begin{frame}{Simulation result: convergence of all structures}
\scriptsize
\begin{columns}[T,onlytextwidth]
\begin{column}{0.31\textwidth}
\begin{block}{What to check}
\begin{itemize}\itemsep2pt
\item Training loss trend.
\item Validation loss trend.
\item Whether series--parallel converges more smoothly than the others.
\end{itemize}
\end{block}
\begin{block}{Interpretation}
Series--parallel is expected to train more stably because it uses measured output history. Parallel is harder because it reuses its own predictions recursively.
\end{block}
\end{column}
\begin{column}{0.65\textwidth}
\centering
\simfig[width=\linewidth,height=0.78\textheight,keepaspectratio]{01_all_structures_cost.png}{Training and validation cost}
\end{column}
\end{columns}
\end{frame}

\begin{frame}{Simulation result: series baseline tracking on unseen 147 mA chirp}
\scriptsize
\begin{columns}[T,onlytextwidth]
\begin{column}{0.31\textwidth}
\begin{block}{Why show series first?}
The series structure is the simplest baseline. It uses only input history and does not receive output history.
\end{block}
\begin{block}{Message}
If series--parallel performs better than this baseline, then measured-output history is helping the identifier.
\end{block}
\end{column}
\begin{column}{0.65\textwidth}
\centering
\simfig[width=\linewidth,height=0.78\textheight,keepaspectratio]{06_series_measured_vs_predicted.png}{Series measured versus predicted}
\end{column}
\end{columns}
\end{frame}

\begin{frame}{Simulation result: series--parallel tracking on unseen 147 mA chirp}
\scriptsize
\begin{columns}[T,onlytextwidth]
\begin{column}{0.31\textwidth}
\begin{block}{Why this is important}
This is Wang's main identification idea during training: use measured past outputs together with past inputs.
\end{block}
\begin{block}{What to look for}
The predicted displacement and force should lie closest to the measured curves. This is the expected advantage of the series--parallel strategy.
\end{block}
\end{column}
\begin{column}{0.65\textwidth}
\centering
\simfig[width=\linewidth,height=0.78\textheight,keepaspectratio]{06_series_parallel_measured_vs_predicted.png}{Series--parallel measured versus predicted}
\end{column}
\end{columns}
\end{frame}

\begin{frame}{Simulation result: parallel free-running tracking on unseen 147 mA chirp}
\scriptsize
\begin{columns}[T,onlytextwidth]
\begin{column}{0.31\textwidth}
\begin{block}{Why parallel is harder}
After initialization, the model feeds back its own previous predictions:
\[
[I,\Delta I,I_{DC},\hat{x},\hat{F}]_{past}
\rightarrow[\hat{x},\hat{F}]_{next}.
\]
\end{block}
\begin{block}{Expected behavior}
Parallel is a stronger test, but its error can accumulate. So it is common to see more drift than in series--parallel one-step prediction.
\end{block}
\end{column}
\begin{column}{0.65\textwidth}
\centering
\simfig[width=\linewidth,height=0.78\textheight,keepaspectratio]{06_parallel_measured_vs_predicted.png}{Parallel measured versus predicted}
\end{column}
\end{columns}
\end{frame}

\begin{frame}{Simulation result: direct structure comparison on the unseen test}
\scriptsize
\begin{columns}[T,onlytextwidth]
\begin{column}{0.48\textwidth}
\begin{block}{Displacement error comparison}
Small non-drifting error means better tracking. This figure compares the three structures directly on displacement.
\end{block}
\centering
\simfig[width=\linewidth,height=0.44\textheight,keepaspectratio]{04_displacement_errors_all_structures.png}{Displacement errors}
\end{column}
\begin{column}{0.48\textwidth}
\begin{block}{Force error comparison}
This figure compares the three structures directly on Lorentz-force prediction error.
\end{block}
\centering
\simfig[width=\linewidth,height=0.44\textheight,keepaspectratio]{04_force_errors_all_structures.png}{Force errors}
\end{column}
\end{columns}
\vspace{0.05cm}
\begin{block}{Key point}
These two plots make the structure comparison much clearer than the old full-record overlay plots, because the errors can be read directly without one diverging curve dominating the vertical axis.
\end{block}
\end{frame}

\begin{frame}{Simulation result: accuracy table and regression comparison}
\scriptsize
\begin{columns}[T,onlytextwidth]
\begin{column}{0.48\textwidth}
\begin{block}{Regression view}
For a strong predictor, the points should cluster near the 45-degree line. This helps visualize prediction quality.
\end{block}
\centering
\simfig[width=\linewidth,height=0.42\textheight,keepaspectratio]{08_all_structures_regression_comparison.png}{Regression comparison}
\end{column}
\begin{column}{0.48\textwidth}
\begin{block}{Accuracy summary}
Reported metrics are RMSE, MAE, $R^2$, and Fit\%. These allow a direct numerical comparison between series, series--parallel, and parallel.
\end{block}
\centering
\simfig[width=\linewidth,height=0.42\textheight,keepaspectratio]{05_all_structures_accuracy_table.png}{Accuracy table}
\end{column}
\end{columns}
\end{frame}

\begin{frame}{Main result from my simulation}
\scriptsize
\begin{columns}[T,onlytextwidth]
\begin{column}{0.48\textwidth}
\begin{block}{What the figures show}
\begin{itemize}\itemsep2pt
\item The \textbf{series} model is a useful baseline.
\item The \textbf{series--parallel} model gives the clearest one-step tracking and usually the best accuracy.
\item The \textbf{parallel} model is the strictest test because it is free-running and can accumulate error.
\end{itemize}
\end{block}
\end{column}
\begin{column}{0.48\textwidth}
\begin{block}{Takeaway about Wang's strategy}
The advantage of the \textbf{series--parallel strategy} is that it uses measured past outputs during training and one-step prediction. This makes learning easier and usually more stable than pure parallel free-running prediction.
\end{block}
\begin{block}{Why keep parallel?}
Parallel is still valuable because it tests whether the learned model can behave like a dynamic simulator when only input and its own past predictions are available.
\end{block}
\end{column}
\end{columns}
\end{frame}

\begin{frame}{Final takeaways}
\begin{block}{Wang (2017)}
The main contribution is nonlinear dynamic system identification using LSTM models and series--parallel/parallel prediction ideas.
\end{block}
\begin{block}{Ogunmolu/Gans (2016)}
The main contribution is evaluating deep dynamic neural-network structures for nonlinear system identification on experimental data.
\end{block}
\begin{block}{My simulation}
The updated simulation trains and compares three structures: series, series--parallel, and parallel. It uses random non-overlapping blocks from 67--127 mA for training/validation and keeps 147 mA completely unseen for the final test.
\end{block}
\end{frame}

\end{document}
